\documentclass{UADYbeamer}
\usepackage[spanish, mexico]{babel}
%---Figuras---TiKZ---
%\usepackage{physics}
\usepackage{amsmath}
\usepackage{tikz}
\usepackage{mathdots}
\usepackage{yhmath}
\usepackage{cancel}
\usepackage{color}
%\usepackage{siunitx}
\usepackage{array}
\usepackage{multirow}
\usepackage{amssymb}
%\usepackage{gensymb}
\usepackage{tabularx}
\usepackage{extarrows}
\usepackage{booktabs}
\usetikzlibrary{fadings}
\usetikzlibrary{patterns}
\usetikzlibrary{shadows.blur}
\usetikzlibrary{shapes}
%--------------------
\title{Título de la plática}
\subtitle{Nombre de la conferencia}
\author[Roberto Balcázar]{Roberto Balcázar \and Otro autor}
\date{\today} % or whatever the date you are presenting in is
\institute[Universidad Autónoma de Yucatán]{Facultad de Matemáticas}
% \copyrightnotice{Published by the American Institute of Aeronautics and Astronautics, Inc., with permission}

\begin{document}

\maketitle

\cutoc

\section{Introducción}

\begin{frame}{Título de la diapositiva}
    \begin{itemize}
        \item Uno
        \item Dos,
        \item $D^{\alpha}\gamma(t)$
        \item Tres
        \item Cuatro
        \item $\begin{pmatrix}
            0&-1\\1&0
        \end{pmatrix}$
        \item Cinco
    \end{itemize}
\end{frame}

\begin{frame}{Otro título}
    \begin{columns}
        \begin{column}{0.6\textwidth}
            \begin{figure}
                \centering
                \input{img/Figura_ejemplo}
                \caption{Figura.}
            \end{figure}
        \end{column}
        \begin{column}{0.4\textwidth}
            \begin{itemize}
                \item Puntos
                \item de
                \item Interés
            \end{itemize}
        \end{column}
    \end{columns}
\end{frame}

\section{Nombre de la sección}

\begin{frame}{Título de la diapositiva}
    \begin{cublock}[Una caja genérica]
        \begin{overlayarea}{\textwidth}{\baselineskip}
            \only<2->{Un texto genérico.}
        \end{overlayarea}
    \end{cublock}
\end{frame}

% Q&A
\begin{frame}{Título}
    \begin{defi}{Un concepto.}{defi_etiqueta}
        Dados $\alpha_1,\ldots,\alpha_n\in (0,1]$ con $\mathbf{\alpha}:=(\alpha_1,\ldots,\alpha_n)$, se define el \textbf{vector fraccional a $\gamma\in C^2((a,b],\mathbb{R}^n)$} como
        $$D^{\mathbf{\alpha}}\gamma(t):=(D^{\alpha_1}\gamma_1(t),\ldots,D^{\alpha_n}\gamma_n(t)).$$
    \end{defi}
\end{frame}

\begin{frame}{Título 2}
    \begin{teo}{Resultado importante}{teo_ejemplo}
        Si $f\in L^1([a,b])$, entonces 
$$D^{\alpha}_{a+}I^{\alpha}_{a+}f(x)=f(x)~~\text{(c.e.t.p.)}.$$
    \end{teo}
\end{frame}
\begin{frame}[standout]
    \Huge\textsc{¡Gracias!}
    
    \vfill
    
    \LARGE\textsc{¿Preguntas?}
\end{frame}

\appendix

\begin{frame}{Diapositivas de respaldo}
    
\end{frame}

\end{document}